\chapter{生态系统也是化学与信息网络}

\begin{kaichang}
下次路过池塘，舀一杯水上来——就用普通的玻璃杯。对着光看，你多半会看到微微发绿，里面漂着几粒几乎看不清的小点。杯底也许沉着一小撮黑泥，闻起来有点腥。

现在猜猜看：这杯看似安静的水里，住着多少“居民”？十个？一百个？再猜猜：去年秋天落进池塘的那些树叶，现在到哪里去了？它们还在池塘里吗？还是“消失”了？

把这两个答案写在纸上。这是全书倒数第二章，我们将把前八十五章攒下的全部家当——原子、化学键、能量、酶、DNA、进化——统统带到这个池塘边，看看它们在行星尺度上合起来演一出什么戏。
\end{kaichang}

\section{落叶并没有消失}

先看能看见的世界。秋天，公园里堆起厚厚的落叶；到了第二年春天，叶子所剩无几。雪化之后你扒开草丛，只能找到一些碎得认不出来的褐色残渣。叶子去哪儿了？

最常见的回答是“烂了”“没了”。但请你用第一章就学过的眼光盯住这个问题：物质由原子组成，原子不会凭空消失（除了核反应，这里没有）。一片叶子里的碳、氢、氧、氮、磷原子，此刻必然身在某个地方。找到它们，就找到了生态系统的入口。

真相是：叶子被\emph{吃掉}了——被细菌、真菌、土壤里的跳虫和蚯蚓，一点点吃掉。这些分解者的细胞里，酶（第 49 章）把叶子中的纤维素、蛋白质拆开，再通过细胞呼吸（第 53 章）把碳架一步步氧化。叶子里的碳原子，一部分变成二氧化碳 \chem{CO_2} 呼回空气，一部分变成分解者身体里的蛋白质，一部分留在土壤里成为腐殖质。氮原子变成铵盐、硝酸盐，被旁边的草根重新吸收，明年又长成新叶。

一个原子都不缺，一个原子也不多。落叶没有消失，它只是\emph{换了一种排列方式}，搬了家。

把这个镜头拉远，你会看到一张网：草被蚱蜢吃，蚱蜢被青蛙吃，青蛙被蛇吃；而蛇、青蛙、蚱蜢、草死后，全都交给分解者。这叫\textbf{食物网}——不是一条笔直的“链”，因为几乎每种生物都有不止一种食物，也不止一个天敌。食物网就是生态系统的“物质流向图”：谁身体里的原子，下一站去谁的身体里。

\section{单行道与环岛路}

食物网里其实流着两样完全不同的东西，它们遵守的规则截然不同。把这两样分开，是理解整个生态系统的关键。

第一样是\textbf{能量}。它几乎总是从同一个入口进来：阳光。生产者（植物、藻类、蓝细菌）通过光合作用把光能转化为化学能，储存在葡萄糖等有机分子里（细节见第 54 章——提醒一句，放出的氧气来自水，不是二氧化碳）。然后能量沿着食物网一级级传递。但每一级传递都要交一笔沉重的“税”：细胞呼吸把大部分化学能变成热，散失到环境中。热一旦散掉，就再也收不回来——它不会重新变成葡萄糖。所以能量在生态系统里走的是\textbf{单行道}：阳光进，热出，一去不复返。整个地球生态系统能持续运转，全靠太阳这个永不关闸的能量源头（以及地球内部的一点余热）。

第二样是\textbf{物质}——碳、氮、磷、硫这些元素的原子。它们走的是\textbf{环岛路}：从空气、水、土壤进入生物体，在食物网里转几手，最后由分解者拆回原形，再被生产者捡起来，开始下一圈。地球从太阳那里几乎只接收能量，几乎不接收物质（每年落下的陨石尘埃相对生物圈可以忽略）。所以生态系统必须精打细算地\emph{回收}每一个原子。

\begin{wuqu}
“能量在食物链里循环利用。”——这是流传很广的说法，但它把能量和物质混为一谈。检验方法很简单：如果能量真的能循环，那么一个完全密封、照不到光的生态瓶应该能永远运转下去。实际上，密封生态瓶一旦放进黑暗柜子，几天到几周内全部居民死光——光合作用停止，没有新能量进来，储存在生物体内的化学能很快以热的形式漏光。反例就在这儿：物质可以在密封瓶里循环很久，能量却必须持续从外界补充。你桌上的台灯、鱼缸里的加热棒、整个生物圈头顶的太阳，都是同一个道理：\textbf{物质靠循环，能量靠补给}。
\end{wuqu}

\begin{qiaoliang}
能量每传一级要交多少“税”？生态学家测量了许多真实的湖泊、草地和森林，发现相邻两级之间，平均只有大约 $10\%$ 的能量被“存进”下一级的身体里（其余在呼吸中变成热、或以未被消化的食物残渣排出）。这个数字很粗糙——不同系统从 $5\%$ 到 $20\%$ 都有——但用来估算威力巨大。

算一笔账：假设一块草地一年通过光合作用固定的能量，折合 $10000$ 单位的化学能。吃草的昆虫能存下约 $10\%\times 10000=1000$ 单位；吃昆虫的鸟约 $100$ 单位；吃鸟的猛禽约 $10$ 单位。也就是说，\textbf{食物网顶端每 1 单位的能量，底层要垫进约 1000 单位}。这直接解释了为什么顶级捕食者总是稀少——一片草原养得起千万棵草、百万只昆虫、几千只鸟，却只养得起几只鹰。也解释了为什么人吃谷物比吃肉“省地”：种谷物喂猪再吃肉，等于多爬一级能量阶梯，同样一块地的产出大打折扣。生态学里这条经验规律叫“百分之十定律”，它的根源不是什么生态学魔法，而是第 27、53 章的热力学：每一次能量转化都有损耗，散掉的热无法回收。
\end{qiaoliang}

用一张图把这两条规则画在一起（这是人为的表示法，方框大小和箭头粗细只是示意）：

\begin{center}
\begin{tikzpicture}[scale=1.0]
\node[draw, rounded corners] (sun) at (0,2.4) {阳光};
\node[draw, rounded corners] (prod) at (0,0.9) {生产者};
\node[draw, rounded corners] (cons) at (3,0.9) {消费者};
\node[draw, rounded corners] (deco) at (6,0.9) {分解者};
\node[draw, rounded corners] (pool) at (3,-0.7) {空气、水、土壤中的元素};
\node (heat) at (7.6,2.4) {热（散失）};
\draw[->, thick] (sun) -- (prod);
\draw[->, thick] (prod) -- (cons);
\draw[->, thick] (cons) -- (deco);
\draw[->, thick] (prod) -- (1.6,0.2) -- (pool);
\draw[->, thick] (deco) -- (pool);
\draw[->, thick] (pool) -- (-0.4,-0.3) -- (prod);
\draw[->, thick, dashed] (prod) -- (1.5,2.4) -- (heat);
\draw[->, thick, dashed] (cons) -- (5,2.4) -- (heat);
\end{tikzpicture}
\end{center}

实线箭头里流着原子（注意那个回到生产者的环），虚线箭头是逐级漏掉的能量。整张图就是本章的标题：生态系统是一台用能量流驱动、靠物质循环维持的机器。

\section{三个循环，一群看不见的工程师}

环岛路不是一条，而是好几条并行的线路，每种元素走自己的路。我们挑对生命最重要的三条：碳、氮、磷。

\textbf{碳循环}你最熟悉。大气中的 \chem{CO_2} 被光合作用固定进有机物，沿食物网传递，再由呼吸和分解放回大气。还有一部分碳走“慢车道”：沉进海底变成石灰岩，或埋入地下变成煤和石油——在那里一锁就是几百万年，直到火山喷发或人类开采才重返大气。碳循环的“大气出入口”由一个化学平衡把守：

\[\chem{CO_2} + \chem{H_2O} \rightleftharpoons \chem{H^+} + \chem{HCO_3^-}\]

海洋吸收的 \chem{CO_2} 越多，平衡越向右移，海水越酸。这条几行字的反应式，后面会看到，正牵动着整个星球的命运。

\textbf{氮循环}的主角不是植物，也不是动物，而是微生物。空气里 $78\%$ 是氮气 \chem{N_2}，但两个氮原子之间的三键极其牢固（第 21 章），植物直接“啃不动”。能把 \chem{N_2} 拆开的，只有固氮细菌——比如住在豆科植物根瘤里的根瘤菌——它们体内的固氮酶在常温常压下完成人类需要高温高压才能做到的事（工业哈伯法，第 44 章）。固氮产物氨 \chem{NH_3} 变成铵 \chem{NH_4^+} 后，硝化细菌把它一步步氧化成硝酸盐 \chem{NO_3^-}：

\[\chem{NH_4^+} \rightarrow \chem{NO_2^-} \rightarrow \chem{NO_3^-}\]

植物吸收硝酸盐，合成蛋白质和 DNA；动物吃植物；排泄物和尸体被分解者拆回铵；最后反硝化细菌在缺氧的泥土里把硝酸盐变回 \chem{N_2}，送回大气，完成闭环。你呼吸的每一口空气里的氮，都可能刚结束一趟穿越细菌和豆科植物的旅行。

\textbf{磷循环}有个关键的不同：它没有“天空航线”。磷几乎没有稳定的气态化合物，只能从岩石风化进入土壤和水，被生物利用，最后随沉积物沉入海底，等亿万年后地壳抬升、重新风化才回来。所以磷是三条循环里最慢、最容易“断供”的一条——许多农田和湖泊里，磷恰恰是限制生物生长的那个短板。

把三条循环摆在一起：

\begin{table}[h]
\centering
\begin{tabular}{llll}
\toprule
循环 & 主要储存库 & 进入生物圈的方式 & 关键微生物角色 \\
\midrule
碳 & 大气、海洋、岩石 & 光合作用固定 \chem{CO_2} & 分解者释放 \chem{CO_2} \\
氮 & 大气（\chem{N_2}） & 固氮作用拆三键 & 固氮、硝化、反硝化细菌 \\
磷 & 岩石、沉积物 & 岩石风化（很慢） & 分解者从尸体中回收 \\
\bottomrule
\end{tabular}
\end{table}

请注意这张表的潜台词：\textbf{没有微生物，循环就断了}。如果细菌、真菌明天集体消失，落叶和尸体将堆积如山，氮被锁死在有机物里，植物得不到氮和磷，食物网从底层崩塌。第 62 章讲过，你的身体本身就是一个由微生物参与的生态系统；现在看到的是放大版：整个生物圈同样靠这些看不见的小家伙当“管道工”。

\section{信息：网络里流动的第三样东西}

能量和物质之外，食物网里还流动着第三样东西：\textbf{信息}，而且它同样是化学的。

蚜虫啃食一株豆苗时，豆苗会向空气中释放特定的挥发性分子；旁边的豆苗“闻”到这些分子，提前启动防御基因，合成让蚜虫难以下咽的物质。玉米被毛虫咬了，释放的信号分子能招来毛虫的天敌寄生蜂——这等于植物在“打电话报警”。真菌用菌丝在森林地下连接多棵树木，树木之间能通过这张“木联网”传递糖、氮和警告信号（细节和争论见第 61 章的边界讨论）。萤火虫用荧光素和荧光素酶的反应（第 53 章提到过的化学发光）闪烁求偶密码；蚂蚁用信息素在地面写下“通往食物”的化学路标。

这些信号的载体，全是我们在这本书里反复见过的分子：萜类、酯类、多肽。接收端则是受体蛋白——和第 58 章讲的细胞信号传导是同一套原理，只是尺度从细胞膜放大到了整片森林。分子发出，扩散，被受体捕获，引发行为改变——这就是信息流。生态系统这张网，既是化学的（物质与能量），也是信息的（信号与响应）。

\begin{shiyan}
\textbf{做一个密封生态瓶，亲眼看看“物质循环、能量补给”。}

材料：带密封盖的透明玻璃罐（约 1～2 升）一个、洗净的细沙或小石子、池塘水或河水一杯、几根水草（金鱼藻、蜈蚣草均可，花鸟市场或水族店有售）、一两只苹果螺（可选）。

步骤：①罐底铺 2 厘米厚的石子；②倒入池塘水至八分满；③放入水草和螺；④盖紧密封盖，放在\emph{明亮但不被太阳直晒}的窗台（直晒会让罐内过热）。⑤贴一张纸条写上日期，之后每周观察一次，持续至少一个月。

观察点：水草是否长出新叶？螺是否活着、在啃食缸壁上的绿膜？内壁上的绿膜（藻类）随时间变多还是变少？把同样配置的第二瓶放进黑暗柜子里做对照，每周对比两瓶的差别。

安全提示：只用池塘水和水草，不放自来水以外的任何化学品；开罐观察后洗手；不要喝瓶里的水；实验结束后把水草和螺倒回来源的池塘或妥善处理，不要倒进下水道或陌生水域（避免物种扩散）。

背后的原理：密封罐只许光进、不许物质进出。水草光合作用制造有机物和氧气；螺吃藻类，呼吸消耗氧气；微生物分解螺的排泄物，释放 \chem{CO_2} 和氮、磷，供水草再用——物质在罐内转圈，能量靠窗外的阳光补给。黑暗对照瓶会告诉你，断掉能量输入后这台小机器多久停转。
\end{shiyan}

\section{多样性：网络的“冗余设计”}

如果你的生态瓶里只放一种水草，一旦它病死，整瓶崩溃；如果有水草、藻类、螺、微生物各司其职，某一种出了状况，别的成员能部分顶替。这就是生态学里一个反复被验证（也反复被细化）的经验规律：\textbf{物种越丰富的生态系统，往往越经得起扰动}。

道理可以用工程的直觉理解。一栋大楼的承重结构如果只有一根柱子，柱子断了楼就塌；有多根柱子，一根坏了其余的分担——工程师管这叫“冗余”。生态系统里，多个物种承担相似功能（比如好几种植物都能固碳，好几种昆虫都能传粉），就等于给网络装了冗余。“生物多样性”这个词听起来像口号，拆开看其实是网络的拓扑性质：节点多、连接密，单点故障就不致命。

不过要小心别把这句话推到极端。多样性不是越多越好、随便什么物种都行——外来物种疯狂增殖反而会拆掉本地网络（比如水葫芦堵死河道）；关键物种（比如传粉者、顶级捕食者）的损失，比同等数量普通物种的损失严重得多。多样性买的是\emph{应对变化的缓冲}，不是永不故障的保险单。

生物多样性还给人类提供一整套“生态系统服务”：植物和微生物净化水源、形成土壤；昆虫为农作物授粉（全球约三分之一的粮食作物依赖动物授粉）；湿地削减洪水；森林固定二氧化碳。这些服务如果改由人工完成——建净水厂、人工授粉、修堤坝——花费是天文数字。四川的果农已经在雇人拿着毛笔给梨花一朵朵授粉，因为本地蜜蜂不够用了：那就是“服务中断”后人工补位的样子，又贵又慢。

\begin{zhengju}
科学家怎么知道“物质在生态系统里循环、且循环会被人为改变”？最直接的证据来自一个大胆的实验：把整片森林“拆开”看看。

1963 年起，美国生态学家博尔曼和莱肯斯在新罕布什尔州的哈伯德布鲁克实验林做了一个经典设计。他们先选中几条相邻的小流域（整片山坡及其溪流），花几年时间测量“收支”：雨水带进来多少钙、钾、硝酸根，溪水带出去多少。结果发现，健康的森林里，大多数营养元素的输出远小于输入——森林把养分\emph{扣}在了自己体内，循环得很紧。

然后是关键一步：1965 年冬天，他们把其中一条流域的树全部砍光，并用除草剂抑制再生，其他流域保持原样作对照。结果溪水里的硝酸根浓度暴涨到砍伐前的几十倍，钙、钾也大量流失——拿掉植被，等于拧开了养分循环的阀门，土壤里的硝化细菌照样生产硝酸盐，却没有根系来吸收，养分被雨水冲个精光。对照流域同期没有这种变化，排除了“那年雨水异常”之类的替代解释。

这个实验的意义不止于“砍树不好”。它第一次用可测量、可对照的方式证明：\textbf{养分循环是生态系统整体的性质}，由植物、土壤、微生物共同维持，拆掉任何一环，循环就变漏。后来同一片森林的长期监测还记录了酸雨的影响和森林几十年的缓慢恢复，成为生态学里“长期实验”的标杆。注意它的方法品格：不靠直觉断言，而是先测基线、设对照、只改一个变量、让数据说话。
\end{zhengju}

\section{当网络被拉扯}

理解了“单行道加环岛路”，就能看懂今天新闻里的生态问题——它们几乎全是这两条规则被人类活动拉扯的结果。

\textbf{富营养化}：农田施的氮肥、磷肥随雨水冲进湖泊，等于给水里藻类空投“军粮”。藻类疯长（水华），死亡后被细菌分解，分解过程大量消耗氧气，鱼和其他生物窒息而死。2007 年太湖蓝藻暴发，让无锡几百万人的自来水发臭，源头正是氮磷输入超过了湖泊循环的处理能力。这不是“自然循环失灵”，而是循环被超量输入\emph{撑爆}了——环岛路进了远超设计流量的车。

\textbf{海洋酸化}：前面那个几行字的平衡式正在全球尺度上移动。人类烧掉化石燃料，大气 \chem{CO_2} 浓度从工业革命前的约 $280\ \mathrm{ppm}$ 升到如今超过 $420\ \mathrm{ppm}$，海洋吸收了其中约四分之一，海水 pH 随之下降。酸化本身对许多生物还不是直接致命的，真正的麻烦在下游：碳酸根 \chem{CO_3^{2-}} 浓度降低，珊瑚、贝类造碳酸钙外壳变得更“费力”（第 39 章沉淀平衡的行星版）。一个平衡被推动，整串依赖它的生命跟着调整。

\textbf{污染物富集}：这是“百分之十定律”的阴暗面。某些物质——比如重金属汞、持久性有机污染物——进入食物网后\emph{不容易}被分解者拆掉，而能量每传一级缩小约十倍，生物体却要把吃进去的食物里的这些物质大体留下。结果是一级一级浓缩：浮游生物身上微不足道，小鱼多一些，大鱼更多，到了食物网顶端的猛禽、大型鱼类和人，浓度可以放大上万倍。上世纪 DDT 让顶级猛禽蛋壳变薄、繁殖失败，正是这个机制。判断一种污染物危不危险，不只看它急性毒性多大，还要看它在食物网里“走不走得掉”。

\textbf{气候变化}：碳循环的慢车道（煤、石油、石灰岩）被人类用两百年时间强行并入了快车道。多出来的 \chem{CO_2} 增强温室效应（第 38 章），地球这个生态系统被迫在新的能量收支下重新找平衡——而生物圈的化学与信息网络，从传粉时间到洋流，全都得跟着改。

把这四件事放在一起看，有一个共同的图景：\textbf{人类活动没有改变任何一条化学规律，只是改变了各种流的大小}——氮流、磷流、碳流、能量流。而生态系统这张网，恰是靠这些流的相对平衡维持的。

\section{模型的边界}

本章的模型什么时候好用，什么时候要小心？

“能量单向、物质循环”在大尺度上非常可靠，但别忘了地球不是绝对封闭：陨石带来少量物质，高层大气每年也逃掉一点氢和氦；人类发射的探测器干脆把地球上的原子永远带出了生物圈。“百分之十定律”是粗平均，具体到某条食物链可能差好几倍，用它算精确账会出错。方框加箭头的循环图把江河湖海简化成几个“库”，真实情况里每个库内部都是不均匀的、变化中的。至于“多样性导致稳定性”，在受控的小尺度实验里证据不错，放大到整个景观尺度，关系会复杂得多——这是第 87 章还要继续面对的开放问题。

\section{回到那杯水}

现在重新看开场那杯池塘水。对着光的那些小点里，有单细胞的藻类在把光子变成化学键；水底的腥臭黑泥里，反硝化细菌在缺氧环境中把硝酸盐一步步变回氮气；杯壁上肉眼难辨的绿膜，是一层微生物组成的“城市”，成员之间用化学信号协调行动。一杯静置的池塘水，只要放在阳光下，可以自我维持很久——因为它是一个微缩的行星：能量从窗口进来，以热的形式离开；碳、氮、磷在杯中循环；信息在分子之间传递。

落叶没有消失，它被拆解、被传递、被重新组装，此刻正以二氧化碳、新叶和细菌身体的形式继续存在。你猜对了吗？

这本书从一个问题开始：物质为什么这样变化？我们从原子出发，走过化学键、反应、能量，走进细胞、DNA、进化，现在站在整个生物圈的高度回望——你会看到同一条因果链从未中断：池塘水里的每一次光合作用和分解，都遵守第 27 章的能量账簿；固氮菌的酶，是第 49 章蛋白质的亲戚；让森林连接成网的信号分子，遵循第 58 章的受体原理；这张网在亿万年间不断变化，靠的是第十篇讲的进化。\textbf{行星尺度上没有新的化学，只有化学的新规模。}

而规模本身会带来新的谜题：这样一张网能被我们预测和管理到什么程度？下一章，也是最后一章，我们去看看科学目前还回答不了的问题——那才是这本书真正的结尾：不是“我们全都知道了”，而是“我们知道到哪里为止”。

\begin{tiaozhan}
试试用地球尺度的数字估算光合作用的份额。太阳辐射到地球大气层顶部的平均功率约为每平方米 $340$ 瓦（考虑了地球是球体、一半时间是夜晚）。地球表面积约 $5.1\times 10^{14}$ 平方米，所以整个行星接收的太阳能功率约 $1.7\times 10^{17}$ 瓦。实测全球光合作用每年固定的化学能约合 $4\times 10^{21}$ 焦耳，除以一年的秒数（约 $3.2\times 10^7$ 秒），平均功率约 $1.2\times 10^{14}$ 瓦。两者相除：光合作用只截获了到达地球太阳能的约 $0.07\%$。就这么一点点，养活了整个生物圈——并且，人类当前的全部能源消耗（约 $2\times 10^{13}$ 瓦）已经达到全球光合固碳功率的约六分之一。数字粗，但量级是实的：我们这个物种的能耗，已经和行星级的“总生产线”可比较了。
\end{tiaozhan}

\section*{练一练}

\begin{sikaoti}
\item 画一张你所在城市一个公园（或校园一角）的物质流图：标出生产者、消费者、分解者各是谁，用箭头画出碳原子可能走的三条路径，并另外用虚线画出能量从进入到离开的全过程。能量和物质的箭头形状有什么根本不同？
\item 一个密封生态瓶放在阳光下能维持数月，放进黑柜子几周就崩溃。用“能量单行道、物质环岛路”解释：密封瓶里到底什么是“循环”的，什么是“漏掉”的？
\item 如果某种除草剂只杀死硝化细菌，而不伤害植物本身，用氮循环的知识预测：几个月后这片农田会发生什么？请画出氮原子流动的变化。
\item 汞在浮游生物体内浓度为 $1$ 个单位，如果每传一级浓缩约 $10$ 倍，吃浮游生物的小鱼、吃小鱼的大鱼、吃大鱼的人，体内浓度大约各是多少？这解释了为什么大型掠食鱼类的食用建议比小型鱼更严格？
\item 太湖蓝藻暴发时，有人提议“多养鱼虾把藻类吃光”。结合能量传递效率和富营养化的成因，分析这个方案为什么通常单独无效。真正需要控制的“阀门”在哪里？
\item 找一则关于“生物多样性保护”的新闻，用本章的“冗余设计”和信息网络的观点写三句话：这个物种或栖息地可能承担着网络里的什么功能？失去它可能牵动哪些连接？
\end{sikaoti}
